\documentclass[11pt]{amsart}
\usepackage[T1]{fontenc}
\usepackage{amsmath,amssymb,amsthm}
\usepackage[margin=1in]{geometry}
\usepackage[colorlinks=true,linkcolor=blue]{hyperref}
\newcommand{\RR}{\mathbb R}
\newcommand{\F}{{\cal F}}
\newcommand{\Fstar}{{\cal F}^{*}}
\newcommand{\dd}{\,dx}
\newtheorem*{editorial}{Editorial note}
\title{Hermite Functions, Fourier Transform, and the Kravchuk System}
\author{Philip}
\date{28 April 2011}
\begin{document}
\maketitle
\begin{center}
\textit{Editable working transcription of FeinNotes11, with marked corrections}
\end{center}

% SOURCE STATUS: Transcribed from the page-by-page readings retained in the
% conversation, including a targeted re-reading of the original scan.
% The original PDF was not available in the current workspace for a fresh
% visual comparison. This is an editorial working transcription, not a
% certified diplomatic reproduction. Section headings and mathematics follow
% the recorded readings; prose and layout have been lightly regularized.
% The recorded handwritten header reads "Philip." and "28/04/11".

\begin{editorial}
This working transcription follows the previously recorded readings of the
three-page notes. The original scan was not available for a fresh visual
comparison during this conversion. Layout and prose have been lightly
regularized. Corrections are recorded in source comments and summarized
at the end; added clarifications are identified explicitly.
\end{editorial}

\section{Hermite functions}
Let $D=d/dx$, and define
\begin{equation}\label{eq:psi}
\psi_n(x)=(x-D)^n e^{-x^2/2}.
\end{equation}
By Heisenberg--Weyl calculus,
\[
x-D=e^{-D^2/2}xe^{D^2/2},\qquad
(x-D)^n=e^{-D^2/2}x^ne^{D^2/2}.
\]
The translation identity is
\[
f(D)e^{ax}=e^{ax}f(D+a).
\]
The generating function is
\begin{align}
G(x,t)&=\sum_{n\ge0}\frac{t^n\psi_n(x)}{n!}
=e^{t(x-D)}e^{-x^2/2}\notag\\
&=e^{-D^2/2}e^{tx}e^{D^2/2}e^{-x^2/2}\notag\\
&=e^{tx}e^{\frac12(D^2-(D+t)^2)}e^{-x^2/2}\notag\\
&=e^{tx-t^2/2}e^{-tD}e^{-x^2/2}\notag\\
&=e^{tx-t^2/2}e^{-(x-t)^2/2}\notag\\
&=e^{2tx-t^2}e^{-x^2/2}.\label{eq:G}
\end{align}
\begin{editorial}
The heat-operator conjugations above are formal Heisenberg--Weyl
identities; separate domain questions for the individual unbounded
exponentials are not being asserted. The final generating function
can alternatively be obtained directly from the Baker--Campbell--Hausdorff
factorization $e^{t(x-D)}=e^{tx}e^{-tD}e^{-t^2/2}$.
\end{editorial}

Define the two-parameter Hermite polynomials by
\begin{equation}\label{eq:Hermite-polynomials}
e^{zx-\tau z^2/2}
=\sum_{n\ge0}\frac{z^n H_n(x,\tau)}{n!}.
\end{equation}
% A targeted re-reading confirmed the factor tau in the exponent.
Then
\[
\psi_n(x)=2^nH_n(x,1/2)e^{-x^2/2}.
\]
For complex $z$, the Gaussian Fourier--Laplace integral is
\[
\frac1{\sqrt{2\pi}}\int_{\RR}e^{zx}e^{-x^2/2}\dd=e^{z^2/2}.
\]

\subsection*{Fourier-transform convention}
% NOTATION CORRECTION requested by the author: write {\cal F}^{*}
% for the positive-exponent transform. Earlier recorded readings called it F.
We write the usual negative-exponent transform as
\begin{equation}\label{eq:Fminus}
(\F f)(\xi)=\frac1{\sqrt{2\pi}}
\int_{\RR}e^{-i\xi x}f(x)\dd,
\end{equation}
and its adjoint (equivalently, its inverse) as
\begin{equation}\label{eq:Fplus}
(\Fstar f)(\xi)=\frac1{\sqrt{2\pi}}
\int_{\RR}e^{i\xi x}f(x)\dd.
\end{equation}
These integrals may first be taken for Schwartz functions and then
extended unitarily to $L^2(\RR,dx)$.
The kernel in \eqref{eq:Fplus} is the complex conjugate of the kernel
in \eqref{eq:Fminus}. For complex-valued $f$, this is not the same as
simply taking the complex conjugate of $\F f$:
\[
\Fstar f=\overline{\F(\overline f)}.
\]

Using \eqref{eq:G}, we obtain
\begin{align}
(\Fstar G)(\xi,t)
&=\frac1{\sqrt{2\pi}}\int_{\RR}
e^{i\xi x}e^{-x^2/2+2tx-t^2}\dd\notag\\
&=e^{-t^2}e^{(i\xi+2t)^2/2}\notag\\
&=e^{-\xi^2/2+2i\xi t+t^2}\label{eq:corrected-mixed-term}\\
&=G(\xi,it)
=\sum_{n\ge0}\frac{(it)^n\psi_n(\xi)}{n!}.\notag
\end{align}
% ORIGINAL RECORDED READING of the intermediate line:
% e^{-xi^2/2+i xi t+t^2}; the coefficient 2 was missing.
% CORRECTED: e^{-xi^2/2+2 i xi t+t^2}.
Comparing coefficients gives
\begin{equation}\label{eq:Fourier-Hermite}
\boxed{\Fstar\psi_n=i^n\psi_n.}
\end{equation}
For the usual transform, correspondingly,
$\F\psi_n=(-i)^n\psi_n$.

\section{$\mathfrak{sl}(2)$ operators}
Consider
\[
H=\frac{x^2-D^2}{2},
\]
the harmonic-oscillator Hamiltonian. The operators
\[
R=\frac{x^2}{2},\qquad L=\frac{D^2}{2},\qquad
\rho=xD+\frac12
\]
satisfy the commutation relations
\begin{equation}\label{eq:commutators}
[L,R]=\rho,\qquad [L,\rho]=2L,\qquad [\rho,R]=2R.
\end{equation}
These differential-operator identities hold on the Schwartz space
$\mathcal S(\RR)$. In particular, $H=R-L$.

By Heisenberg--Weyl algebra,
\[
(x+D)(x-D)=x^2-D^2+1,
\qquad 2H=(x+D)(x-D)-1.
\]
Now
\[
(x-D)G(x,t)=\sum_{n\ge0}\frac{t^n\psi_{n+1}(x)}{n!},
\]
% An earlier reading rendered the generating parameter here as l.
% It has been regularized to the same t used throughout G(x,t).
and
\begin{align*}
(x+D)G(x,t)
&=(x+D)e^{2tx-t^2-x^2/2}\\
&=2tG(x,t)
=\sum_{n\ge1}\frac{2n\,t^n\psi_{n-1}(x)}{n!}.
\end{align*}
Thus
\[
(x+D)\psi_n=2n\psi_{n-1},\qquad
(x-D)\psi_n=\psi_{n+1},
\]
where the lowering term vanishes at $n=0$. Consequently,
\begin{align*}
2H\psi_n
&=\bigl[(x+D)(x-D)-1\bigr]\psi_n\\
&=\bigl[(2n+2)-1\bigr]\psi_n
=(2n+1)\psi_n.
\end{align*}
Therefore
\[
e^{\pi iH/2}\psi_n
=e^{\pi in/2+\pi i/4}\psi_n
=i^n\sqrt{i}\,\psi_n,
\qquad \sqrt{i}=e^{i\pi/4}.
\]
Together with \eqref{eq:Fourier-Hermite}, this gives the main identity
\begin{equation}\label{eq:main-on-basis}
\boxed{e^{\pi iH/2}\psi_n=\sqrt{i}\,\Fstar\psi_n.}
\end{equation}

\subsection*{Hilbert-space formulation}
% ORIGINAL RECORDED SENTENCE:
% "This extends to L^2(R,e^{-x^2/2} dx) since the functions psi_n
% are dense in that Hilbert space."
% EDITORIAL CORRECTION: the standard unitary Fourier/oscillator identity
% belongs to unweighted L^2(R,dx), for these Gaussian wavefunctions.
The normalized Hermite functions
\[
h_n(x)=\frac{\psi_n(x)}{\sqrt{2^n n!\sqrt\pi}}
\]
form an orthonormal basis of $L^2(\RR,dx)$. Taking the self-adjoint
oscillator realization of $H$, both sides of
\eqref{eq:main-on-basis} define bounded unitary operators on this
space. Equality on the Hermite basis therefore implies
\begin{equation}\label{eq:main-operator}
\boxed{e^{\pi iH/2}=\sqrt{i}\,\Fstar,\qquad
\Fstar=e^{\frac{\pi i}{2}(H-\frac12)}.}
\end{equation}
Equivalently,
\[
\F=e^{-\frac{\pi i}{2}(H-\frac12)}.
\]
A Gaussian-weighted polynomial realization is also possible, but
requires conjugating the operators; it is not the same unitary
Fourier transform with a different measure inserted.

\medskip
\noindent\textit{Remark from the notes.}
The relation between the $\mathfrak{sl}(2)$ algebra and Fourier
transform is well known in optics.
% This is retained as a source remark, not a newly verified literature claim.

\section{Connection with Kravchuk polynomials}
The notes state:
\begin{quote}
The Kravchuk polynomials are basis states for irreducible
representations of $\mathfrak{sl}(2)$. They are eigenfunctions of
$H$ in a different representation.
\end{quote}
\begin{editorial}
This is retained as a programmatic representation-theoretic statement.
Here $H=(x^2-D^2)/2$ denotes the differential oscillator Hamiltonian.
A corresponding finite-dimensional generator must be specified in its
own representation; the quoted sentence is not an assertion that these
finite matrices are the same operator as the differential $H$.
No finite-dimensional convolution intertwiner is established here.
\end{editorial}

\section{Convolution in the Kravchuk system}
The final page poses the question:
\begin{quote}
Find an intertwining operator that relates convolution/Fourier
transform with K-convolution/Kravchuk transform. Does it exist???
\end{quote}
This remains an open question in these preliminary notes.

\appendix
\section{Editorial changes and points for review}
\begin{enumerate}
\item \textbf{Fourier notation.} Following the author's clarification,
$\Fstar={\cal F}^{*}$ denotes the positive-exponent transform.
Its Hermite eigenvalues are $i^n$. The negative-exponent transform
$\F$ has eigenvalues $(-i)^n$.
\item \textbf{Mixed-term typo.} The intermediate Gaussian exponent
is $-\xi^2/2+2i\xi t+t^2$, not
$-\xi^2/2+i\xi t+t^2$.
\item \textbf{Hermite parameter.} The generating function for
$H_n(x,\tau)$ contains $-\tau z^2/2$. This agrees with the targeted
re-reading recorded earlier and with the specialization
$2^nH_n(x,1/2)$.
\item \textbf{Generating parameter.} A previously transcribed $l^n$
in the raising identity has been regularized to $t^n$; this may have
been a reading ambiguity rather than a source typo.
\item \textbf{Hilbert space.} The Fourier/oscillator unitary identity
is stated on $L^2(\RR,dx)$, not on the Gaussian-weighted space in
the earlier recorded sentence. The original reading is preserved
in a source comment.
\item \textbf{Scope.} The formal heat-operator calculation, the
finite-versus-continuum representation issue, and the unsolved
convolution question have explicit editorial qualifications.
\item \textbf{Source comparison.} A fresh page-by-page comparison
with the original scan remains desirable before treating this as
an exact transcription. No new theorem or literature attribution
has been added on the author's behalf.
\end{enumerate}
\end{document}
